\documentclass[../../main.tex]{subfiles}
\begin{document}

\begin{flushright}
  \footnotesize\textit{【自動文字起こし・要確認】}
\end{flushright}

% Source: handwritten.pdf 第2問．
{\bf [解答]}
$n$秒のあいだに，2頂点$A,B$にしかおとずれない確率を$a_n$，
3頂点$A,B,C$にしかおとずれない確率を$b_n$と定める．
$a_n$について，$P$は$A\to B\to A\to B\to\cdots$の移動をくり返すので，
\begin{align}a_n=\left(\dfrac{1}{3}\right)^n.\label{eq:1}
\end{align}
対称性から$(A,C)$，$(A,D)$にしかおとずれない確率も$a_n$である．
$b_n$について，常に3頂点のうち上記のごとく2頂点のみに移動をつづけ，
かつ一度は$A,B,C$全てにおとずれることから
\begin{align}
  b_n=\left(\dfrac{2}{3}\right)^n-2a_n
  =\left(\dfrac{2}{3}\right)^n-2\left(\dfrac{1}{3}\right)^n.
  \label{eq:2}
\end{align}
対称性から，$(A,B,D)$，$(A,C,D)$のみにおとずれる確率も$b_n$である．
以上から，求める確率を$c_n$として，排反から
\begin{align}
  c_n=1-3a_n-3b_n
  =1-3\left\{\left(\dfrac{2}{3}\right)^n-\left(\dfrac{1}{3}\right)^n\right\}.
\end{align}

\end{document}
